%
\subsection*{Transformada unilateral y existencia}
\label{CH06-B01}
La transformada de Laplace es
$\mathcal L\{f\}(s)=\int_0^\infty e^{-st}f(t)\,dt$.
La continuidad a trozos en cada intervalo finito y una cota
$|f(t)|\le M e^{at}$ para $t$ suficientemente grande son condiciones
suficientes de convergencia para $s>a$. Hay que comprobar por separado
el origen y la cola: el factor exponencial no corrige una singularidad
no integrable en $t=0$.
\[
\begin{array}{c|c|c}
 f(t)&\mathcal L\{f\}(s)&\text{condición suficiente}\\\hline
 t^n&n!/s^{n+1}&s>0,\ n=0,1,\ldots\\
 e^{at}&1/(s-a)&s>a\\
 \cos\omega t&s/(s^2+\omega^2)&s>0\\
 \sin\omega t&\omega/(s^2+\omega^2)&s>0
\end{array}
\]
En las dos últimas filas, $\omega>0$. Las condiciones indicadas se
refieren a $s$, no al tiempo. Recuérdense integrales impropias y
exponenciales.

%
\subsection*{Propiedades y reloj de la señal}
\label{CH06-B02}
Si $F=\mathcal L\{f\}$, la transformación es lineal y
$\mathcal L\{e^{at}f(t)\}(s)=F(s-a)$. Para $a\ge0$, el escalón
$H(t-a)$ vale cero antes de $a$ y uno después; su valor aislado en $a$
no altera la transformada. La traslación temporal se escribe
\[
 \mathcal L\{H(t-a)f(t-a)\}(s)=e^{-as}F(s).
\]
Para transformar $H(t-a)g(t)$, póngase primero
$f(\tau)=g(a+\tau)$: $\tau$ mide el tiempo transcurrido desde la
activación. Reexpresar sólo una parte de la señal cambia el problema.

Si $b>0$, $\mathcal L\{f(bt)\}(s)=F(s/b)/b$. Para una señal de periodo
$T>0$, integrable y continua a trozos en un periodo,
\[
 \mathcal L\{f\}(s)=\frac{\int_0^T e^{-st}f(t)\,dt}{1-e^{-sT}},
 \qquad s>0.
\]
Una señal que alcanza una meseta no es por ello periódica.
Consulta: solución de \exref{M2-E0378}.

%
\subsection*{Transformar los datos e invertir}
\label{CH06-B03}
Con regularidad y crecimiento suficientes para integrar por partes,
\[
 \mathcal L\{y'\}=sY-y(0),\qquad
 \mathcal L\{y''\}=s^2Y-sy(0)-y'(0).
\]
En orden $n$, se restan todos los términos
$s^{n-1-j}y^{(j)}(0)$, $j=0,\ldots,n-1$. Sustitúyanse los datos con
sus signos antes de despejar $Y$. Las fracciones parciales y completar
cuadrados permiten reconocer la tabla; los factores $e^{-as}$ requieren
una traslación temporal.

Sepárese la respuesta debida a los datos de la debida a la entrada y
compruébese después la función recuperada en la ecuación y en el PVI.
Una solución puede existir para todo $t\ge0$ y crecer exponencialmente;
su transformada sólo convergerá para ciertos valores de $s$. Esta región
de convergencia no es un intervalo temporal. Consulta: solución de
\exref{M2-E0122}.

%
\subsection*{Cambios de operación y continuidad}
\label{CH06-B04}
Una señal que cambia de ley en $a$ puede escribirse como ley anterior
más $H(t-a)$ por la diferencia de leyes, expresada con el reloj adecuado.
Para varios cambios se suman las correcciones. Si una fuerza es finita
y continua a trozos, en el oscilador con masa no nula se conservan
posición y velocidad al cambiar la fuerza; la aceleración puede saltar.
Estas condiciones también enlazan las soluciones obtenidas por tramos.

Laplace resulta útil para reunir entradas activadas; resolver por tramos
puede ser más breve si hay pocos eventos. En ambos casos deben
comprobarse los empalmes, los datos y el cierre de un ciclo cuando se
pida. Que $x'(a)=0$ no demuestra que $a$ sea la primera parada, ni que
el cuerpo permanezca detenido. Apagar la fuerza tampoco añade
amortiguamiento. Consultas: soluciones de \exref{M2-E0043} y
\exref{M2-E0314}.
